\begin{abstract}
\NEW{Atmospheric turbulence displaces the apparent positions of sources in a single image by a spatially correlated, curl-free field
whose variance grows as the exposure time shortens, which makes it the dominant source of astrometric error for the 30\,s visits
of the Vera C. Rubin Observatory. In the 28\,518 visits of Data Preview 2 (DP2), we measure an $E$-mode variance of
${\sim}340\,\mathrm{mas}^2$ in the astrometric residuals of LSSTCam, ten times the HSC value and three times the DES value.
DP2 corrected this field with the Gaussian process of \citet{2021A&A...650A..81L}, a sheared von K\'arm\'an kernel fit to the
measured two-point correlation function, which removes ${\sim}90\%$ of the turbulence variance but is slow, memory-hungry, and
fails on a fraction of the visits. We present TURBO-GP (TUrbulence Removal By fourier-Optimized Gaussian Process), which is 
inpsired by the work of \citet{2025AJ....170..361G} and pushed it further, by the uses of the measured two-point correlation function 
directly as the kernel, with nothing to fit, and solves the Gaussian process on a regular grid of the focal plane by fast Fourier 
transform and conjugate gradient, with
positive semi-definiteness imposed once in Fourier space. On DP2, TURBO-GP reduces the error rate of the turbulence-fitting task
from 5\% to 1\%, the total wall time from ${\sim}2.7$ years to ${\sim}2$ weeks, and the typical memory per job from ${\sim}7$\,GB
to ${\sim}1$--$2$\,GB, while leaving less residual correlation; it also allows the training sample to grow from 10\,000 to 50\,000
sources per visit, which halves the residual $E$-mode variance on validation sources. The corrected residuals are a clean probe
of the atmosphere, whose $E$-mode correlates with the PSF FWHM, and, averaged in sensor coordinates, they reveal the static
distortions of LSSTCam, of ${\sim}4$\,mas on ITL and 1--2\,mas on E2V sensors.}
\end{abstract}
